% Avsnitt 13.8, ur lectures/L13/appendix/b_exercises.md.

\section{Övningar}
\label{c13:sec:exercises}

\begin{aside}
De här övningarna befäster kapitlet. Konstruera din egen \file{crc15.vhd} utifrån
\secref{c13:sec:iface} till \ref{c13:sec:worked}; den utdelade testbänken, körd enligt
bilaga~\ref{app:sim}, är kontrollen av att den beter sig korrekt.
\end{aside}

\exercise{\code{crc15}}{VHDL}
\label{c13:ex:crc15}
Skriv din egen \file{crc15.vhd} utifrån specifikationen i det här kapitlet. Verifiera den sedan:

\begin{shell}
cd controller
ghdl -a --std=93 can_def.vhd crc15.vhd crc15_tb.vhd
ghdl -e --std=93 crc15_tb
ghdl -r --std=93 crc15_tb --assert-level=error
\end{shell}

Om en kontroll fallerar, läs assertion-meddelandet (bilaga~\ref{app:sim}).

\exercise{CRC för hand, sedan i simulering}{Simulering}
\label{c13:ex:byhand}
\xpart{a} Räkna för hand ut, med \secref{c13:sec:recursion}:s rekursion, vilket CRC-15-resultat som
blir följden av att mata in 4-bitarssekvensen \code{1010} (MSB först) i en nyss nollställd motor.
Visa uträkningen en bit i taget (återkoppling, skift, villkorlig XOR).

\xpart{b} Bekräfta ditt handräknade svar med GHDL: skriv en kort testbänk (eller anpassa strukturen i
\file{crc15_tb.vhd}) som nollställer \code{crc15}, matar in \code{1010} och rapporterar det
resulterande \code{crc}-värdet, antingen via en \code{assert} eller via en
\code{report ... severity note;} som du läser direkt.

\xpart{c} Utgången \code{valid} hos \code{crc15} läser \code{'1'} direkt efter reset, innan någon
riktig data alls har matats in. Förklara varför det inte är en bugg, med hänvisning till vad
\code{can_controller} (\chapref{c17:ch}) faktiskt gör med \code{valid}, och när.

\exercise{Felet som en summerad checksumma missar}{Simulering}
\label{c13:ex:missed}
\Secref{c10:sec:framing} påstod att en CRC ”fångar en betydligt bredare, matematiskt välförstådd
klass av fel” än den enkla bytevisa summa som dess eget ramformat använder. Den här övningen gör det
konkret genom att köra en korruption som summan bevisligen inte kan se genom motorn du just byggt.

Ta nyttolasten \code{0x3C, 0x91} och korruptionen \code{0x34, 0x99}: bit 3 nollställd i den första
byten, bit 3 satt i den andra, så att den ena byten förlorar exakt det den andra vinner.

\xpart{a} Räkna ut den bytevisa checksumman för båda nyttolasterna och bekräfta att de är
\textbf{identiska}: båda blir \code{0xCD}. En enkel summa registrerar bara totalen, så en mottagare
som kontrollerar den skulle godta den korrupta ramen som hel.

\xpart{b} Mata nu in alla 16 bitarna av den \emph{ursprungliga} nyttolasten i en nyss nollställd
\code{crc15} (MSB först, första byten först) och notera det resulterande \code{crc}-värdet. Nollställ
sedan och gör samma sak med den \emph{korrupta} nyttolasten. Stämmer de två CRC-värdena överens?
Använd simulering; en \code{report} av båda värdena räcker.

\xpart{c} Egenskapen att generering och kontroll är samma operation (\secref{c13:sec:twojobs}) säger
att en ram vars egna CRC-bitar matas tillbaka in återför registret till noll. Mata in den
\textbf{korrupta} nyttolasten följd av den \textbf{ursprungliga} nyttolastens CRC-bitar, alltså
precis det en mottagare skulle se om den här korruptionen inträffade under transporten. Når registret
noll? Går \code{valid} hög? Förklara vad mottagaren drar för slutsats.

\xpart{d} Säg med en mening vilken egenskap hos CRC-rekursionen som gör att den märker en förändring
som en summa inte kan, med hänvisning till vad varje bit gör med registret i
\secref{c13:sec:recursion}.
